\section{Discussion}

The strongest evidence in Phase 3 is negative but useful. It shows that a reference-conditioned region-prompt route cannot be assumed simply because privileged GT boxes work. The frozen DINOv2-L/local-peaks evidence can find high-precision fragments, yet the fragments do not produce region proposals that beat matched random controls with sufficient query-level consistency. The gap between privileged P4 and automatic R1 therefore becomes an automatic-proposal gap, not a SAM refinement result.

This distinction protects the scientific claim. If we had skipped R1 and run SAM on weak proposals, a later Dice failure would conflate proposal quality with the segmenter's ability to refine a region. Conversely, if a few qualitative examples looked good, they would not justify method claims under the pre-registered win-rate and matched-control thresholds. R1 forces the pipeline to prove that the automatic prompt input is already meaningful before spending SAM calls.

The result also reframes AdaPromptSet. The original budget-constrained prompt-set selection story remains too strong for the available evidence. A more accurate direction is reference-conditioned prompt representation analysis: determining when a point, box, mask, direct matching mask, or deterministic region is justified by frozen evidence. In this GlaS setting, the answer is that the tested automatic region proposals are insufficient.

One practical consequence is that the next legal experiment cannot be a quiet threshold expansion or per-query strict-validation tweak. It would require a new pre-registered rescue protocol using only allowed development data or a new external data source, with explicit data-manifest and method gates. Without that, the correct final state is \texttt{STOP\_GLAS\_REGION\_PROPOSAL\_FAILURE}.
